\section*{Chapter \ref{ch:g7:volume-mass-density} --- Volume, Mass, Density}

\begin{solution}{exo:g7:volume-mass-density:1}
$\rho = m/V$: $\rho$ the \omterm{def:g7:volume-mass-density:formula}{density}, $m$ the sample's \omterm{def:g3:measuring-mass:mass}{mass}, $V$ its
\omterm{def:g6:mass-and-volume:volume}{volume}. With \omterm{def:g3:measuring-mass:units}{grams} and \omterm{def:g6:mass-and-volume:volume}{cubic centimetres}, $\rho$ comes out in
\unit{g/cm^3}.
\end{solution}

\begin{solution}{exo:g7:volume-mass-density:2}
$\rho = 270 \div 100 = \qty{2.7}{g/cm^3}$: aluminium.
\end{solution}

\begin{solution}{exo:g7:volume-mass-density:3}
$m = 9.0 \times 40 = \qty{360}{g}$ of copper.
$V = 1930 \div 19.3 = \qty{100}{cm^3}$ of gold --- the crown
case's own number.
\end{solution}

\begin{solution}{exo:g7:volume-mass-density:4}
Taring removes the cylinder's own \omterm{def:g3:measuring-mass:mass}{mass} so the balance reports the
\omterm{def:g3:solids-liquids-gases:liquid}{liquid} alone --- the honest-zero rule. The formula is then fed
the balance's \omterm{def:g3:measuring-mass:mass}{mass} $m$ and the \omterm{met:g6:mass-and-volume:cylinder}{meniscus} reading $V$.
\end{solution}

\begin{solution}{exo:g7:volume-mass-density:5}
$\rho = 70 \div 50 = \qty{1.4}{g/cm^3}$ --- honey-like. Ice at
$0.9$ is far less dense: it \omterm{def:g1:floating-and-sinking:words}{floats} high in it.
\end{solution}

\begin{solution}{exo:g7:volume-mass-density:6}
Both \omterm{def:g6:measurement-in-science:measuring}{units} scale together: a \omterm{def:g6:mass-and-volume:volume}{litre} is $1000$ \unit{cm^3} and a
\omterm{def:g3:measuring-mass:units}{kilogram} is $1000$ \omterm{def:g3:measuring-mass:units}{grams}, so the two thousands cancel --- one
number serves both. Mixing \omterm{def:g3:measuring-mass:units}{grams} with \omterm{def:g6:mass-and-volume:volume}{litres} skips that
cancellation and delivers a number a thousandfold wrong.
\end{solution}

\begin{solution}{exo:g7:volume-mass-density:7}
$\rho = 58 \div 5.0 = \qty{11.6}{g/cm^3}$ --- near lead's
$11.3$, nowhere near gold's $19.3$. Verdict: not pure gold;
likely a lead-hearted impostor in gold clothing.
\end{solution}

\begin{solution}{exo:g7:volume-mass-density:8}
$100 \times 100 \times 100 = 1000000$ \unit{cm^3} --- a million.
At \qty{1}{g} each, that is \qty{1000000}{g} $=$ \qty{1000}{kg}:
a full tonne per cubic \omterm{def:g2:measuring-length:units}{metre} of water --- and why a large
waterbed is furniture for the ground floor.
\end{solution}

\begin{solution}{exo:g7:volume-mass-density:9}
$V = 260 - 200 = \qty{60}{cm^3}$;
$\rho = 420 \div 60 = \qty{7.0}{g/cm^3}$ --- between aluminium
($2.7$) and iron ($7.9$), closest below iron. Between-ness
whispers: a blend or a plated impostor, certainly not silver's
$10.5$.
\end{solution}

\begin{solution}{exo:g7:volume-mass-density:10}
\omterm{def:g3:solids-liquids-gases:liquid}{Liquid} \omterm{def:g3:measuring-mass:mass}{mass}: $1010 - 380 = \qty{630}{g}$; \omterm{def:g6:mass-and-volume:volume}{volume}:
$\qty{75}{cL} = \qty{750}{cm^3}$; so $\rho = 630 \div 750 =
\qty{0.84}{g/cm^3}$ --- an oil.
\end{solution}

\begin{solution}{exo:g7:volume-mass-density:11}
Parcel \omterm{def:g7:volume-mass-density:formula}{density}: $1200 \div 4000 = \qty{0.3}{g/cm^3}$ --- far
below water's $1.0$: the buoy \omterm{def:g1:floating-and-sinking:words}{floats} high. Aluminium's $2.7$
describes only the metal skin; the parcel is mostly enclosed \omterm{def:g2:air-around-us:air}{air},
and it is the parcel --- total \omterm{def:g3:measuring-mass:mass}{mass} over total \omterm{def:g6:mass-and-volume:volume}{volume} --- that
faces the floating rule.
\end{solution}

\begin{solution}{exo:g7:volume-mass-density:12}
Protocol: weigh the chain ($m$); \omterm{def:g6:measurement-in-science:measuring}{measure} its \omterm{def:g6:mass-and-volume:volume}{volume} by the rise
method, fully submerged, no bubbles ($V$); compute $\rho = m/V$;
purity expects about \qty{9.0}{g/cm^3}. Honest doubts: hollow
links trap \omterm{def:g2:air-around-us:air}{air} and swell $V$, faking a low \omterm{def:g7:volume-mass-density:formula}{density}; and a cunning
blend (or a plated core) can land near $9.0$ while containing no
pure copper at all. Tightening test: repeat the \omterm{def:g6:mass-and-volume:volume}{volume}
measurement after flooding the links (shake out bubbles), and add
an independent signature --- for instance the \omterm{def:g1:magnets:magnet}{magnet} (a steel
core betrays itself instantly) or, in a workshop, a measured
\omterm{def:g6:water-states:changes}{melting} behavior.
\end{solution}

\begin{solution}{pb:g7:volume-mass-density:1}
\textbf{1.} $5400 \div 2000 = \qty{2.7}{g/cm^3}$: aluminium.
\textbf{2.} $1800 \div 200 = \qty{9.0}{g/cm^3}$: copper.
\textbf{3.} $4520 \div 400 = \qty{11.3}{g/cm^3}$: lead --- dense,
soft, and poisonous to handle carelessly: gloves.
\textbf{4.} $8100 \div 1500 = \qty{5.4}{g/cm^3}$ --- between
aluminium ($2.7$) and iron ($7.9$): a mixed sack of the two,
roughly half and half by \omterm{def:g6:mass-and-volume:volume}{volume}.
\textbf{5.} $V = 540 \div 2.7 = \qty{200}{cm^3}$.
\textbf{6.} One cube: $m = 9.0 \times 25 = \qty{225}{g}$. The
coil's \qty{1800}{g} yield $1800 \div 225 = 8$ full cubes.
\textbf{7.} $\qty{20}{kg} = \qty{20000}{g}$; $20000 \div 225 =
88.9$: $88$ cubes.
\textbf{8.} $m = 7.9 \times 3000 = \qty{23700}{g} \approx
\qty{23.7}{kg}$.
\textbf{9.} $3860 \div 200 = \qty{19.3}{g/cm^3}$ --- a perfect
match for gold.
\textbf{10.} Tungsten --- \omterm{def:g7:volume-mass-density:formula}{density} $19.3$, the same signature to
the decimal. The \omterm{def:g7:volume-mass-density:formula}{density} test identifies candidates; it cannot
distinguish substances that happen to share a signature.
\textbf{11.} Neither instrument failed: \omterm{def:g3:measuring-mass:mass}{mass} and \omterm{def:g6:mass-and-volume:volume}{volume} are
correct, and so is the division. The test honestly established
that the ingot is \emph{either} gold or something of gold's exact
\omterm{def:g7:volume-mass-density:formula}{density} --- it narrowed the suspects to two.
\textbf{12.} Defensible choices: the struck bar's ring and feel
(tungsten is far harder --- a file or hardness test on a hidden
corner tells them apart quickly), or an accepted expert \omterm{def:g6:measurement-in-science:measuring}{measure}
of how the bar conducts \omterm{def:g5:heat-and-insulation:heat}{heat} or current (gold is among the best
\omterm{def:g4:conductors-insulators:conductor}{conductors}, tungsten far behind) --- gentler than any furnace,
and decisive together with \omterm{def:g7:volume-mass-density:formula}{density}.
\textbf{13.} For example: ``\omterm{def:g7:volume-mass-density:formula}{Density} can prove a substance is
\emph{not} what it claims. It can only suggest what it is ---
signatures narrow the suspects. And alone it can never convict:
identity wants two independent witnesses.''
\end{solution}
